% 例 2：正八边形每内角 135°、每外角 45°
\documentclass[tikz,border=4pt]{standalone}
\usepackage{ctex}
\usepackage{amsmath}
\usetikzlibrary{calc}
\begin{document}
\begin{tikzpicture}[scale=1.0, thick]
  \def\R{2.2}
  \foreach \i in {0,...,7} {
    \coordinate (P\i) at ({22.5 + \i*45}:\R);
  }
  \draw (P0)
    \foreach \i in {1,...,7} { -- (P\i) } -- cycle;
  \foreach \i in {0,...,7} \filldraw (P\i) circle (1.2pt);
  % 把右下顶点 P7 的相邻两边向外延长一小段：示意外角 45°
  \coordinate (Q) at ($(P0)!-0.35!(P7)$); % P0 反向延长
  \draw[gray, dashed, thick] (P0) -- (Q);
  % 内角标注 135°（在最上 P1 处简化标注）
  \node[font=\footnotesize] at (0,0) {正八边形};
  \node[font=\footnotesize, above right] at (P0) {$135°$};
  \node[font=\footnotesize, red] at ($(P0)+(0.5,-0.55)$) {$45°$};
\end{tikzpicture}
\end{document}
